\section{Lean formalization: definitions, obligations, and object proofs}
\label{sec:formalization}

\subsection{The exact status of the supplied Lean files}

The packet contains two Lean source outlines:
\begin{itemize}
\item \code{lean/FourSquares.lean}: total executable definitions and proposed
proof scripts assembling the theorem from six explicit arithmetic contracts.
It imports Lean's standard library and the \code{omega} tactic, without
importing Mathlib's four-square theorem.
\item \code{lean/ProveItTarget.lean}: the exact first-order iPRA goal in
ProveIt's syntax, a stronger goal parameterized by four unary primitive
recursive symbols, and a proposed object-level Euler-identity proof fragment.
This file requires the full pinned ProveIt workspace.
\end{itemize}

\textbf{Neither file was elaborated or compiled in this environment: a Lean
compiler was not available.} No completed Lean four-square proof or final
$\mathsf{PRAAx}$ derivation is claimed. The files expose the work still to
be done as named local obligations; they contain no admitted theorem
declarations. Absence of an admission is not a compilation receipt, and
a conditional theorem is not an unconditional proof of its conclusion.

The preceding sections prove the mathematical content of the obligations.
The separate formalization ledger, \code{OBLIGATIONS.md}, decomposes the
remaining work into 31 concrete steps, including the final syntax-level
translation.

\subsection{Definitions and their correspondence to the proof}

The host Lean development uses an explicit record with four integer
coordinates, and two records for the loops. The complete definitions are in
the file, with this correspondence:

\begin{longtable}{@{}L{0.42\textwidth}L{0.52\textwidth}@{}}
\toprule
Mathematical object & Lean identifier\\
\midrule
\endhead
Finite conjunction and bounded first hit & \code{allLt}, \code{firstLt}\\
Primality bit and first divisor & \code{primeB}, \code{IsPrime}, \code{leastDivisor}\\
Integer four-tuple, norm, product, conjugate &
\code{Z4}, \code{Z4.normSq}, \code{Z4.mul}, \code{Z4.conj}\\
Centered residues & \code{centered}, \code{centered4}\\
Three pairings and exact halving & \code{halfPairs}, \code{evenHalf}\\
Initial multiple of a prime & \code{seedState}\\
Transition~\eqref{eq:step} & \code{descentStep}\\
State assertion~\eqref{eq:valid}, rank & \code{PrimeState}, \code{descentRank}\\
Primitive recursive iteration & \code{iterate}\\
Prime witness~\eqref{eq:primeprog} & \code{primeRep}\\
Outer factor transition & \code{factorStep}\\
Outer norm invariant and rank & \code{FactorInvariant}, \code{factorRank}\\
Signed result and natural result & \code{signedFourSquares}, \code{fourSquares}\\
Unfinished arithmetic interface & \code{ArithmeticContracts}\\
Conditional correctness and existence & \code{fourSquares\_correct},
\code{lagrange\_from\_contracts}\\
\bottomrule
\end{longtable}

\par\noindent
\begin{minipage}{\linewidth}
For example, the recursion is structurally on its explicit clock:
\begin{lstlisting}
def iterate (step : A -> A) : Nat -> A -> A
  | 0, s     => s
  | k + 1, s => step (iterate step k s)
\end{lstlisting}
\end{minipage}\par
Here \code{A} is a host-level type parameter. The actual iPRA instances
use natural state codes and fixed first-order primitive recursive operations.
The type-polymorphic presentation does not imply that iPRA itself has
higher-order objects.

\subsection{The six local contracts}

The file's \code{ArithmeticContracts} structure contains exactly the
following mathematical claims:
\begin{enumerate}
\item \textbf{Seed correctness:} a prime input yields a valid seed state.
This is Proposition~\ref{prop:seed}, including finite-search correctness.
\item \textbf{Descent preservation:} $I_p(s)$ implies $I_p(\Step(s))$.
This is Proposition~\ref{prop:step}.
\item \textbf{Descent rank decrease:}
$R(\Step(s))\le R(s)\monus1$ for $R=m\monus1$ on valid states.
The terminal, even, and odd branches were proved separately.
\item \textbf{Factor-loop preservation:} assuming the proved prime witness
property, the outer transition preserves
$q\ge1$ and $q\Nrm{A}=n$.
\item \textbf{Factor-loop rank decrease:} on this invariant,
$(q'\monus1)\le(q\monus1)\monus1$.
\item \textbf{Naturalization:} replacing each integer by its natural
absolute value preserves the sum of squares.
\end{enumerate}

The source proposes generic preservation and rank-bound proofs, then uses
these six fields to assemble
\begin{lstlisting}
theorem fourSquares_correct (C : ArithmeticContracts) (n : Nat) :
    (fourSquares n).normSq = n

theorem lagrange_from_contracts (C : ArithmeticContracts) (n : Nat) :
    exists a b c d : Nat, n = a*a+b*b+c*c+d*d
\end{lstlisting}
The actual Lean files use standard Lean existential notation. The argument
\code{C} is an explicit hypothesis record, and no instance of it is asserted.
To finish the host formalization, prove the listed elementary lemmas,
construct that record, and specialize these statements to it.

Useful local proof methods are elementary rewriting for quotient equations,
\code{omega} for linear inequalities and truncated subtraction, and a
polynomial normalization tactic for the two Euler identities and the
halving identity. Nonlinear inequalities such as
$(m-1)^2<m^2$ should be reduced to explicit polynomial expansions plus
positive-factor reasoning; they should not be silently given to an
unrestricted high-level theorem.

\subsection{The required distinction: host Lean truth versus iPRA derivability}

A theorem of type
\[
 \forall n:\NN,\;\exists a,b,c,d:\NN,\quad n=a^2+b^2+c^2+d^2
\]
is a host Lean theorem about natural numbers. On its own, it does not
state that iPRA proves the sentence. A complete machine formalization of
the foundational claim must construct an object derivation of type
\[
 ([]:\operatorname{List}(\mathsf{Fm}\ \mathsf{PTm}\ 0))
 \ \vdash[\mathsf{PRAAx}]\ \mathsf{fourSquareSentence}.
\]
This type is explicitly defined as \code{FourSquareObjectGoal} in the
second source file; no inhabitant is postulated.

The innermost matrix has five variables. Inside
$\forall n\,\exists a\,\exists b\,\exists c\,\exists d$, ProveIt's
de Bruijn indices are
\[
 0=d,\quad1=c,\quad2=b,\quad3=a,\quad4=n.
\]
Thus the matrix is
\[
 v_4\doteq(v_3v_3+v_2v_2)+v_1v_1+v_0v_0,
\]
with the association fixed in \code{natSquareSum}. The four existential
constructors are wrapped inside one universal constructor. This specifies
the target without any ambiguity about which witness each binder denotes.

For the stronger result, the source defines a family parameter
$W:\operatorname{Fin}4\to\mathsf{PR}\,1$ and the closed object sentence
\[
 \forall n\;\left(n=\sum_{i<4}W_i(n)W_i(n)\right).
\]
The remaining construction is to compile the four actual programs into
four \code{PR 1} terms and produce the derivation of that identity.
Existential introductions then derive the ordinary sentence.

Host-level generic induction lemmas may quantify over arbitrary predicates.
Their use here is legitimate for the iPRA audit only because the actual
instantiations are the coded quantifier-free predicates in
Section~\ref{sec:audit}. A machine proof must record these QF certificates
and invoke \code{PRA.qfInd} or an already proved restricted induction
interface. Merely observing that Lean accepted a host induction proof
would not establish this restriction.

\subsection{Concrete integration with the inspected ProveIt snapshot}

The source review is pinned to commit
\begin{center}
\texttt{3a521d0265941f9e9d4ad450818365a78a6c9a73}.
\end{center}
At that snapshot the root toolchain specifies Lean \code{v4.32.0};
the dependency documentation records Mathlib commit
\texttt{81a5d257c8e410db227a6665ed08f64fea08e997}~\cite{proveit}.
The following actual source interfaces supply relevant object derivations:

\begin{longtable}{@{}L{0.34\textwidth}L{0.60\textwidth}@{}}
\toprule
Source module & Integration use\\
\midrule
\endhead
\code{Core.Definitions} & \code{PTm.toPR}, \code{recDef},
\code{rec\_zero}, \code{rec\_succ} compile the fixed arithmetic programs
and expose their recursion equations.\\
\code{Core.Search} & \code{Search.search\_spec} and
\code{Search.search\_found} give bounded-search success, minimality,
and failure information. Its search returns the bound on failure.\\
\code{Core.Pigeonhole} & \code{Pigeonhole.pigeonhole} supplies explicit
colliding indices for a coded list of $K+1$ values below $K$. Its import
closure is substantial; the report also proves the simple counting argument
directly.\\
\code{Core.EuclidLemma} & \code{euclid\_lemma\_at} supplies the prime
divisor disjunction as an object derivation.\\
\code{AnalyticNumberTheory.Primes} & \code{lpf\_dvd} and
\code{lpf\_prime} justify a bounded least-prime-factor search.\\
\code{Core.Pairing} & \code{pairT}, \code{fst\_pair},
\code{snd\_pair}, and \code{simF} support state coding and simultaneous
primitive recursion.\\
\code{Core.Integers} & \code{ZT} represents differences of natural terms;
\code{ZT.derive} and \code{ipra\_ring} can produce the fixed integer
polynomial identities as object proofs.\\
\code{Core.ClockedDescent} & \code{clocked\_at} and \code{clocked}
give the QF refutational descent principle. The witness-producing absorbing
loop here should instead use its displayed rank invariant directly.\\
\bottomrule
\end{longtable}

These are source-inspected interfaces, not newly recompiled dependencies.
The copied reference files retain their source identifiers and the
repository's MIT-0 license. They are a reference subset, not a transitive
import closure. Build the object target in the full pinned repository.

There is an important encoding distinction. ProveIt's \code{ZT} is a pair
$(u,v)$ representing $u-v$, with no canonicality condition. Equality of
$(u,v)$ and $(r,s)$ is the natural equation $u+s=r+v$. The canonical
integer code in Appendix~\ref{app:coding} has an explicit PR bridge to
this representation: normalize a difference pair to
$(u\monus v,v\monus u)$, then encode its sign and nonzero magnitude.
The polynomial identity can be proved in \code{ZT} first and transported
through this bridge. Thus no assumption about quotient-type integer
objects is needed inside iPRA.

\code{Core.ClockedDescent} has a complete explicit proof in the inspected
source, and its introducing commit reports a successful build and axiom
audit. This packet did not independently recompile it, and no independent
CI check result was retrieved~\cite{clocked}. Moreover, that interface
forbids a nonempty QF state predicate closed under strict descent; it does
not already supply the tuple-producing absorbing iteration of this report.
Keeping that difference explicit prevents a refutational theorem from
being mistaken for a finished witness constructor.

A sensible new module family is
\code{IPRADerivations/DiophantineNumberTheory/FourSquares/}, with modules
\code{Identity}, \code{Seed}, \code{Halving}, \code{OddDescent},
\code{PrimeWitness}, \code{Assembly}, and \code{Statement}.
Their dependency order is the order of this report. The six local
contracts and the QF invariants specify the precise mathematical boundary
between them.

\subsection{Relation to the existing ordinary Lean theorem}

Mathlib already provides \code{Nat.sum\_four\_squares}, together with the
Euler identity and a prime case. The inspected source uses
\code{classical}, \code{Nat.findX}, and induction on multiplicative
factorization~\cite{mathlib}. This is useful as a comparison implementation.
It is not used as a premise here: the explicit bounded witnesses and QF
induction audit are what establish the requested weak-theory result.

\begin{samepage}
The inspected ProveIt file is
\begin{quote}\small
\path{Analysis/FabiusFunction/Lean/FabiusFunction/FourSquareTheorem.lean}
\end{quote}
It uses Mathlib's theorem to obtain an ordinary existence corollary. Its
source is distinct from a $\mathsf{PRAAx}$ derivation~\cite{proveit}.
The source inventory found no four-square-named module in the inspected
iPRA derivation tree; this is a filename observation, not a claim that
no related material occurs anywhere in the repository.
\end{samepage}
